% !TEX root = ../main.tex
%----------------------------------------------------------
% appendix-ds-theory.tex — Connections to Dempster–Shafer Theory
% Moved from §8 (Discussion) to appendix in Round 3 revision.
%----------------------------------------------------------

\section{Connections to Dempster--Shafer Theory}%
\label{sec:ds_connections}

The possibility--necessity duality has a well-known structural
counterpart in Dempster--Shafer (DS) theory
\citep{Yager2005-at}. In DS terms, the possibility measure
$\poss(A)$ corresponds to the plausibility function $\mathrm{Pl}(A)$:
both represent an upper bound on the probability of the event $A$.
Conditional necessity $\Nc(A)$ plays the role of the belief function
$\mathrm{Bel}(A)$: both represent a lower bound. The interval
$[\Nc(A),\,\poss(A)]$ used throughout this paper is therefore
structurally equivalent to the DS interval
$[\mathrm{Bel}(A),\,\mathrm{Pl}(A)]$, and the verification metrics
defined in Section~\ref{sec:native_verification}
can be read as special cases of
interval-probability verification.

A key difference concerns the nature of the residual uncertainty.
Possibilistic ignorance $\ign = 1 - \max(\pi)$ quantifies a
\emph{coverage gap}: the fraction of plausibility that the forecasting
system left unassigned to any outcome. In DS theory, uncertainty
instead arises from \emph{evidential conflict} among multiple
information sources, and is distributed across focal elements of
varying specificity. The scalar $\ign$ therefore has no direct DS
analogue; it is a distinctly possibilistic quantity. The
possibility-to-probability bridge
(Section~\ref{sec:poss_prob_bridge}) is inspired by the pignistic
transformation of \citet{Smets1990-nf} but extends it to
subnormal distributions: the classical pignistic transform assumes a
normal (i.e., $\max(\pi)=1$) mass function, whereas the bridge
explicitly allocates the subnormality residual $\ign$ to a dedicated
ignorance outcome before redistributing the remaining mass.

The structural parallels suggest a natural avenue for generalisation.
The verification framework developed here could be extended to
belief-function forecasts by replacing $\poss$ and $\Nc$ with
$\mathrm{Pl}$ and $\mathrm{Bel}$ throughout. The five-number scorecard
extends naturally: $\alpha^{*}$ becomes the plausibility
of the observed outcome, and $\eta$ measures the diffuseness of the
plausibility function. We note
this connection as a direction for future work, as a full treatment
would require addressing the richer combinatorial structure of general
mass functions.
